\part{Composition Keyboards}

\chapter{The Linotype}

The Linotype shows that ``typing'' has never named a single kind of work. An operator at the keyboard does not produce a page. The keyboard releases brass matrices that assemble into a line of type, which is then cast as one piece.

\section{Mechanism}

Ottmar Mergenthaler's Linotype was developed in the 1880s and entered newspaper production in the middle of that decade. \needs{dates and first installations} Pressing a key releases a matrix from a magazine; matrices and spacebands accumulate to fill a line; the line is cast in metal and the matrices are returned to their channels for reuse.

\section{A Keyboard Not Built on QWERTY}

The Linotype keyboard is not arranged like a typewriter's. Its lowercase keys follow an order derived from letter frequency, the source of the familiar nonsense phrase ``etaoin shrdlu'' that appears in printing lore when an operator runs a finger down the first columns to fill a line. \needs{verify the arrangement and the origin of the phrase} A training curriculum for this keyboard could not reuse the typewriter's home-row lessons, and the operator's speed was governed as much by the machine's cycle as by the hand.

\section{Consequences for Tutoring}

Two points matter for later chapters. First, the layout is evidence that the arrangement of keys is a design variable chosen for a purpose. Second, the operator's error is different in kind: a wrong key produces a wrong matrix that must be discarded or fixed in the line, so errors have material cost in metal and time.

\chapter{Monotype and Separated Keying}

The Monotype system, developed by Tolbert Lanston, divides composition into two machines. \needs{dates} A keyboard perforates a paper ribbon, and a separate caster reads the ribbon and casts individual characters.

\section{The Ribbon as Intermediate Record}

Separating keying from casting inserts a record between the hand and the metal. The ribbon can be inspected, corrected, and re-run. It also carries the information needed to justify a line, so the operator's task includes tracking the space remaining in the line.

\section{Why It Matters Here}

The ribbon is an early instance of a persistent record of keystrokes that exists apart from the finished product. The paper page of Chapter~8 and the practice ledger of Chapter~27 belong to the same family, and Monotype shows that the family predates computing.

\chapter{Keypunches and Teletypes}

The keyboards of data processing detach the visible mark from the mechanism.

\section{Keypunches}

A keypunch translates each keystroke into a pattern of holes in a card. The operator's product is a card that no reader can interpret without a machine, though many keypunches also printed the character along the top. Eighty-column cards became the dominant format in data processing. \needs{dates and vendors}

\section{Teleprinters}

Teleprinters and terminals transmit codes, and the mark is produced elsewhere. Later character encodings standardized what each keypress means. \needs{ASCII dates and the relationship between teleprinter codes and computer keyboards}

\section{Emitting a Code}

With the computer keyboard the typist no longer strikes a letter onto paper but emits a code that is interpreted somewhere else. That detachment is what makes a typing tutor possible in the sense of this book: the same computer that receives the trained gesture can observe, score, and reshape it.

\chapter*{Study Questions, Part III}
\addcontentsline{toc}{chapter}{Study Questions, Part III}
\begin{enumerate}
\item What kinds of error are costly on a composition keyboard but not on a typewriter?
\item A perforated ribbon is a persistent record of keystrokes. Name one respect in which it is like the typed page of the paper tutor and one in which it differs.
\item Why does the detachment of mark from mechanism make an automatic tutor possible?
\end{enumerate}
